\documentclass[11pt]{article}
\usepackage{fontspec}
\setmainfont{Times New Roman}
\usepackage[a4paper,top=2.3cm,bottom=2.2cm,left=2.3cm,right=2.3cm]{geometry}
\usepackage{microtype}
\usepackage{enumitem}
\usepackage{titlesec}
\usepackage{booktabs}
\usepackage{tabularx}
\usepackage{ragged2e}
\usepackage[hidelinks]{hyperref}

\titleformat{\section}{\normalfont\large\bfseries}{\thesection.}{0.5em}{}
\titlespacing{\section}{0pt}{12pt}{4pt}
\setlist[itemize]{leftmargin=1.3em,topsep=2pt,itemsep=2pt,parsep=0pt}
\setlength{\parindent}{0pt}
\setlength{\parskip}{5pt}
\linespread{1.04}
\frenchspacing
\newcommand{\wc}[1]{\mbox{WC-#1}}
\newcolumntype{L}{>{\RaggedRight\arraybackslash}X}

\begin{document}
\justifying

\begin{center}
{\LARGE\bfseries WeConnect 2.0 --- Retest Report}\\[4pt]
{\large Independent verification of the reported fixes}\\[10pt]
\end{center}

\begin{tabular}{@{}ll@{}}
\textbf{Prepared for} & Lighthouse Communities Foundation \\
\textbf{Date} & 20 October 2026 \\
\textbf{Source retested} & current \texttt{dev} of all 18 repositories, plus live \texttt{api-dev} \\
\end{tabular}

\vspace{3pt}\hrule\vspace{2pt}

You reported 32 Fixed, 20 Not a Bug, 7 Partially Fixed and 23 Pending. We retested by re-reading the
current source and re-running the automated checks against it, rather than reading the status labels
back to you. Where a fix is real we say so; where it is partial or absent we show the evidence. We do
not yet have the per-bug list of which identifiers you placed in each bucket --- once we do, we will
map this one to one and respond to the twenty marked ``Not a Bug''. The identifiers below are ours,
from the register and tracker already with you.

\section{Summary}
\begin{tabularx}{\textwidth}{@{}lcL@{}}
\toprule
\textbf{Outcome} & \textbf{Count} & \textbf{Findings} \\
\midrule
Verified fixed & 3 & \wc{011}, \wc{130}, \wc{095} \\
Partially fixed & 6 & \wc{131}, \wc{129}, \wc{068}, \wc{073}, \wc{133}, \wc{134} \\
Still present & 11 & \wc{084}, \wc{088}, \wc{079}, \wc{105}, \wc{128}, \wc{013}, \wc{004}, \wc{121}, \wc{112}, \wc{023}, \wc{116} \\
Runtime re-confirm & 9 & \wc{048}, \wc{084}, \wc{095}, \wc{035}, \wc{122}, \wc{125}, \wc{126}, \wc{078}, \wc{115} \\
\bottomrule
\end{tabularx}

\section{Verified fixed}
\begin{itemize}
  \item \textbf{\wc{011}} --- the client-side auth guard is no longer hardcoded \texttt{true}; it is
        derived from the credentials. Confirmed in every frontend app.
  \item \textbf{\wc{130}} --- the stored-XSS sink is gone. No \texttt{dangerouslySetInnerHTML} or
        \texttt{innerHTML} assignment remains anywhere in the frontend.
  \item \textbf{\wc{095}} (S1, cross-centre invoice) --- the invoice API now has a centre-scope
        check. It exempts \texttt{System~Manager}, but the live accounts no longer hold that role, so
        the check now applies to them. Source-verified; the runtime write is being re-confirmed on the
        isolated stack (section 5).
\end{itemize}

\section{Partially fixed}
\begin{itemize}
  \item \textbf{\wc{131}} --- the reset OTP is now drawn with a cryptographic generator
        (\texttt{secrets.randbelow}). It is still a four-digit code, so the space is 9000 values.
  \item \textbf{\wc{129}} --- the key is no longer hardcoded in a source file, but the live key is
        still committed in \texttt{.env}, \texttt{.env.prod} and \texttt{.env.uat}. It needs rotating
        and purging from history, since the repository was public.
  \item \textbf{\wc{068}} --- the \texttt{youth\_skilling} API spec endpoints now require
        authentication; the \texttt{batch\_management} spec is still served to anonymous callers.
  \item \textbf{\wc{073}} --- the live accounts no longer hold \texttt{System~Manager}: a real
        improvement. But the code still bypasses permissions at the same number of sites, so the
        model is now enforced inconsistently rather than centrally (see the Remediation Plan).
  \item \textbf{\wc{133}} --- a React error boundary now exists in two of nine apps; the other seven
        still blank out on any component fault.
  \item \textbf{\wc{134}} --- 204 images without \texttt{alt}, down from 264.
\end{itemize}

\section{Still present}
Verified in the current source or live, on 20 October:
\begin{itemize}
  \item \textbf{\wc{084}} (S1) --- \texttt{close\_batch} still has no centre or permission check.
  \item \textbf{\wc{088}} --- the \texttt{id = json.loads(id)} typo is still at both sites.
  \item \textbf{\wc{079}}, \textbf{\wc{105}} --- the permission bypass is unchanged: 322
        \texttt{ignore\_permissions} sites and 66 \texttt{System~Manager} references in backend code.
  \item \textbf{\wc{128}} --- the password-reset flow still logs the OTP to the browser console. We
        had recorded this as fixed on 8 October; on re-reading the current code it is not, and we are
        correcting that here.
  \item \textbf{\wc{013}}, \textbf{\wc{004}} --- the API secret is still kept in \texttt{localStorage}
        (15 sites) and 19 \texttt{.env} files are still committed.
  \item \textbf{\wc{121}}, \textbf{\wc{112}}, \textbf{\wc{023}}, \textbf{\wc{116}} --- CORS reflects
        any origin, no security headers are sent, tracebacks still appear in error bodies, and the
        session is seven days. All still live on \texttt{api-dev}. If these are fixed in a branch not
        yet deployed there, tell us which environment to retest.
\end{itemize}

\section{Runtime re-confirmation (isolated stack)}
Rebuilt from the current \texttt{dev} source and seeded with synthetic data. The two critical
write breaches were tested end to end.
\begin{itemize}
  \item \textbf{\wc{084}} (S1) --- \textbf{still present, confirmed at runtime.} A Skilling Partner
        account permitted only for East and West centres closed a batch owned by North Centre: the
        API answered \emph{``Batch closed successfully.''} \texttt{close\_batch} applies no centre
        check.
  \item \textbf{\wc{095}} (S1) --- \textbf{fixed for ordinary accounts, confirmed at runtime.} A
        non-\texttt{System~Manager} account permitted only for East and West, raising a North
        invoice, was refused with HTTP 403, \emph{``You do not have permission to raise this
        invoice.''} The breach now reproduces only for an account that holds \texttt{System~Manager};
        the live accounts no longer do, so they are protected. The fix is therefore conditional on
        account-role configuration rather than enforced for every caller --- if any privileged
        account keeps \texttt{System~Manager}, the breach returns for it.
\end{itemize}
The remaining write-side items (\wc{048}, \wc{035}, \wc{122}, \wc{125}, \wc{126}, \wc{078},
\wc{115}) share the same authorisation root cause; their status stands on the source and earlier
runtime evidence and was not re-executed in this session.

\section{What we need from you}
\begin{itemize}
  \item The per-bug list behind your 32 Fixed / 20 Not a Bug, so we map this report one to one and
        respond to each ``Not a Bug''.
  \item Which environment carries the configuration fixes, if \texttt{api-dev} is not it.
  \item A Skilling Partner account that is actually linked to a partner, so SP-scoped paths can be
        retested (the one issued returns ``not linked to any Skilling Partner'').
\end{itemize}
The grouped plan for everything still open is in the Remediation Plan accompanying this report.
\end{document}
